\documentclass[12pt, a4paper]{article}

\usepackage[utf8]{inputenc}
\usepackage[T1]{fontenc}
\usepackage{lmodern}
\usepackage[margin=2.5cm]{geometry}
\usepackage[round]{natbib}
\usepackage{hyperref}
\usepackage{booktabs}
\usepackage{url}
\usepackage{setspace}

\doublespacing

\bibliographystyle{apalike}

% ---------------------------------------------------------------
\title{The Verification Gap: Why AI-Augmented Academic Writing\\
       Needs Reporting Guidelines Beyond Citation Checking}

\author{[Author Name]\textsuperscript{1}\\[4pt]
        \small\textsuperscript{1}[Affiliation]\\
        \small [email]}

\date{}
% ---------------------------------------------------------------

\begin{document}
\maketitle

% ===============================================================
% ABSTRACT
% ===============================================================
\begin{abstract}
Large language models are increasingly used to assist with academic
writing, from literature synthesis to argument structuring.  While
citation hallucination---the fabrication of plausible-sounding
references---has received considerable attention, it represents only one
failure mode among several.  AI writing assistants may also inflate the
confidence of speculative claims, expand arguments beyond available
evidence, and produce fluent prose that obscures gaps in reasoning.
Existing solutions, including citation verification tools and
retrieval-augmented generation, operate at the model or tool level; to
our knowledge, no process-level verification infrastructure addresses
the writing workflow itself.  This gap is compounded by a structural
absence in reporting guidelines: the EQUATOR Network maintains
approximately 500 guidelines for empirical health research, yet
essentially no equivalent frameworks exist for non-empirical paper
types---theoretical, design science, perspective, and methodological
contributions.  Drawing on Gregor's (2006) theory-type taxonomy, we
identify two theory types (analytic and design) that are fundamentally
non-empirical and for which argument quality, rather than factual
accuracy, is the primary verification challenge.  We propose a typed
verification model that distinguishes three unit types---claims
(source-verifiable facts), arguments (Toulmin-verified reasoning
chains), and propositions (Whetten-verified recommendations with
boundary conditions)---each requiring distinct verification procedures.
Preliminary evidence from three retrospective audits suggests that typed
verification may reveal actionable issues invisible to conventional
approaches, including false failures caused by applying the wrong
verification procedure to the wrong unit type.  We propose that the
scholarly community begin developing verification infrastructure for
non-empirical papers analogous to what EQUATOR provides for empirical
research.

\medskip\noindent
\textbf{Keywords:} academic writing, artificial intelligence, reporting
guidelines, verification, argumentation, non-empirical research
\end{abstract}

\clearpage

% ===============================================================
% SECTION 1 — THE PROBLEM
% Registry entries: S1-1, S1-2, S1-3, S1-4
% ===============================================================
\section{Introduction: The Problem}

The rapid adoption of large language models (LLMs) as writing assistants
is transforming how academic papers are produced.  These tools can
synthesise literature, structure arguments, generate prose, and format
references with remarkable fluency.  Yet this fluency introduces risks
that are distinct from---and in some ways more subtle than---the
challenges of traditional academic writing.  Where a human author's
limitations are typically visible in the text (awkward phrasing, gaps in
coverage, inconsistent formatting), an AI-assisted draft can be polished
and comprehensive while harbouring verification failures that are
invisible to casual inspection.

% S1-1: Citation hallucination (CLAIM, P0, EMERGING)
The most widely discussed risk is citation hallucination.  LLMs can
invent plausible-sounding papers, complete with fabricated authors,
journal names, and digital object identifiers.  Unlike factual errors
that contradict known evidence, hallucinated citations create a
distinctive category of falsehood: references that cannot be falsified
through domain knowledge alone because they refer to works that do not
exist.  Citation hallucination has received considerable attention, and
tools designed to detect it are emerging.

% S1-2: Confidence inflation (CLAIM, P0, EMERGING)
However, citation hallucination may be only the most visible failure
mode.  A second, potentially more insidious risk is confidence
inflation: AI writing assistants may state speculative claims with the
same certainty as well-established findings, using language such as
``demonstrates'' or ``confirms'' where ``suggests'' or ``may indicate''
would be more appropriate.  In academic writing, the distinction between
``demonstrates'' and ``suggests'' carries substantial epistemic
weight---it signals to readers whether a claim rests on strong evidence
or preliminary findings.  When AI-generated prose conflates these
registers, the resulting text can misrepresent the state of evidence
without containing any factually incorrect statements.  \citet{liang2024}
found that LLM-generated feedback on research papers catches surface
issues but struggles with the kind of deep argument analysis that would
detect such confidence miscalibration.  A retrospective audit of a
technology paper under review at an IEEE journal identified six of
twenty-two verifiable entries where the manuscript language was more
confident than the underlying evidence warranted.

% S1-3: Scope creep (CLAIM, P1, EMERGING)
A third risk is scope creep.  Without structural constraints such as
section specifications and word budgets, AI-assisted drafts may expand
arguments beyond what the available evidence supports, introduce
tangential sections, and exceed format requirements.  While scope
management is a challenge in all academic writing, AI tools can amplify
it by generating additional content with minimal effort, making
expansion the path of least resistance.

% S1-4: Existing solutions at model/tool level (CLAIM, P0, EMERGING)
These failure modes share a common characteristic: they cannot be
detected by checking whether citations are real.  Citation verification
tools such as RefChecker and scite.ai, along with model-level approaches
including retrieval-augmented generation and grounded generation, address
the factual accuracy dimension of AI-assisted writing.  These are
valuable contributions, but they operate at the model or tool
level---they improve what the AI generates or verify whether specific
references exist.  To our knowledge, no process-level verification
infrastructure addresses the writing workflow itself: the systematic
registration, tracking, and confidence calibration of claims throughout
the drafting process.

The distinction between tool-level and process-level verification
matters because the failure modes described above are not primarily
failures of the AI model.  They are failures of the writing
process---failures that occur when there is no infrastructure to register
claims at appropriate confidence levels, calibrate language to evidence
strength, or verify that arguments are logically sound rather than merely
well-sourced.  This paper asks what such process-level infrastructure
might look like and whether existing scholarly frameworks provide a
foundation for building it.


% ===============================================================
% SECTION 2 — THE LANDSCAPE GAP
% Registry entries: S2-1, S2-2, S2-3, S2-4
% ===============================================================
\section{The Landscape Gap}

% S2-1: EQUATOR ~500 guidelines (CLAIM, P0, EMERGING)
To assess the current state of verification infrastructure for academic
writing, it is useful to examine the most comprehensive effort to date:
the EQUATOR Network.  EQUATOR (Enhancing the QUAlity and Transparency Of
health Research) maintains a database of approximately 500 reporting
guidelines for health research \citep{equator}.  These
guidelines---including CONSORT for randomised controlled trials, STROBE
for observational studies, and PRISMA for systematic reviews
\citep{page2021}---provide structured checklists that authors complete
during manuscript preparation.  The GoodReports tool
\citep{struthers2021} automates guideline selection based on study type,
further reducing the barrier to adoption.  The impact of these guidelines
on the quality and transparency of empirical research reporting is widely
recognised; they represent the most mature infrastructure for structured
verification in academic publishing.

% S2-2: No guidelines for non-empirical papers (CLAIM, P0, EMERGING)
However, an examination of the EQUATOR database reveals a structural
gap.  Virtually all of its approximately 500 guidelines address empirical
research: studies that collect, analyse, and report data.  Systematic
reviews, covered by PRISMA \citep{page2021}, represent a partial
exception---they are structured and methodical but synthesise rather than
generate primary data.  Beyond this partial exception, to our knowledge,
no EQUATOR or equivalent guidelines exist for theoretical papers, design
science contributions, perspective articles, or methodological papers
that do not report empirical findings.

% S2-3: Gregor's theory types (CLAIM, P1, EMERGING)
This absence is not accidental.  It reflects a fundamental difference in
what needs to be verified across paper types.  \citet{gregor2006}
identified five theory types in information systems research based on
four goals: analysis, explanation, prediction, and prescription.  Two of
these types---Type~I (analytic theories, including taxonomies and
frameworks) and Type~V (design and action theories, including design
principles and methods)---are fundamentally non-empirical.  Their
contributions are not data points that can be verified against
measurements or observations; they are conceptual structures whose value
lies in their clarity, coherence, and utility.  The remaining types
(Type~II: explanatory, Type~III: predictive, Type~IV: explanatory and
predictive) align more naturally with empirical verification approaches,
as they make claims about observable phenomena.

While Gregor's taxonomy was developed for information systems, the
underlying distinction applies broadly.  Perspective articles,
commentaries, theoretical syntheses, and methodological proposals all
share the characteristic that their primary contribution is an argument
or a framework rather than empirical evidence.  These paper types appear
across virtually every academic discipline, from the humanities and
social sciences to engineering and medicine.

% S2-4: Argument quality is primary challenge (CLAIM, P1, EMERGING)
For these non-empirical paper types, the central verification challenge
may not be ``do the data support the claims?'' but rather ``is the
argument logically sound, are the premises established, and are the
boundary conditions specified?'' \citep{toulmin2003, whetten1989}.  This
is a fundamentally different verification task from the one that EQUATOR
guidelines are designed to support.

This suggests that the EQUATOR gap is structural rather than a mere
oversight.  Empirical papers need verification of research design, data
collection, analysis, and reporting---domains for which structured
checklists are well-suited.  Non-empirical papers need verification of
argument quality: warrant validity, reasoning coherence, confidence
calibration, and boundary specification.  These are different
verification tasks that may require different infrastructure.  The
question is whether established frameworks from argumentation theory and
theory-building methodology can provide a foundation for this
infrastructure.


% ===============================================================
% SECTION 3 — TYPED VERIFICATION: THE PROPOSAL
% Registry entries: S3-1, S3-2, S3-3, S3-4
% ===============================================================
\section{Typed Verification: A Proposal}

We propose that verification infrastructure for non-empirical papers
should distinguish three unit types, each requiring a distinct
verification procedure (Table~\ref{tab:types}).

\begin{table}[htbp]
\centering
\caption{Three verification unit types with per-type procedures.}
\label{tab:types}
\small
\begin{tabular}{@{}llll@{}}
\toprule
\textbf{Type} & \textbf{Definition} & \textbf{Verification question}
  & \textbf{Framework} \\
\midrule
Claim       & Source-verifiable fact
  & Does the source say this?  & Citation checking \\
Argument    & Evidence + inference $\to$ conclusion
  & Is the reasoning valid?    & Toulmin (1958/2003) \\
Proposition & Recommendation with boundaries
  & Are conditions specified?  & Whetten (1989) \\
\bottomrule
\end{tabular}
\end{table}

The first and most familiar type is the \textbf{claim}: a factual
statement that can be verified against one or more sources.  Claims are
the default unit in citation-based verification---``Study~X found~Y'' is
a claim whose accuracy can be checked by reading Study~X.  Most existing
verification tools, including citation checkers and fact-verification
systems, operate on claims.  For claims, the verification question is
straightforward: does the cited source exist, and does it say what the
manuscript attributes to it?

% S3-1: Toulmin for arguments (CLAIM, P1, EMERGING)
The second type is the \textbf{argument}: a reasoning chain that
combines evidence with inference to reach a conclusion.  Arguments
appear throughout academic writing, particularly in discussion and
conclusion sections, yet they resist claim-level verification because
their validity depends not on any single source but on the logical
relationship between premises and conclusion.  \citet{toulmin2003}
provides a well-established framework for analysing argument structure,
decomposing any argument into six components: claim (the conclusion
asserted), grounds (the evidence offered), warrant (the inferential
bridge connecting grounds to claim), backing (support for the warrant),
qualifier (the degree of certainty), and rebuttal (conditions under
which the argument does not hold).  This framework may provide a basis
for operationalising argument verification in academic writing: an
argument can be evaluated by checking whether its grounds are
established, its warrant is explicit and valid for the target audience,
and its qualifier is calibrated to the strength of the evidence.  Recent
work by \citet{gupta2024} found that LLMs can extract Toulmin-style
argument structures through zero-shot prompting, suggesting that
automated support for argument-level verification may be technically
feasible.

% S3-2: Whetten for propositions (CLAIM, P1, EMERGING)
The third type is the \textbf{proposition}: a novel recommendation or
contribution that specifies what should be done and under what
conditions.  Propositions are the defining unit of many non-empirical
papers---they are what the paper contributes to the field.
\citet{whetten1989} provides a framework for evaluating theoretical
contributions through four components: What (the constructs or factors
involved), How (the relationships between them), Why (the underlying
logic or reasoning), and Who/Where/When (the boundary conditions
specifying where the contribution applies).  Propositions can be
evaluated by checking whether all key constructs are defined, whether the
reasoning is explicit, and whether boundary conditions are
stated---without requiring empirical data.

% S3-3: Confidence-to-language mapping (CLAIM, P1, EMERGING)
These three types may benefit from a shared confidence calibration
system.  We propose mapping confidence tiers to prescribed language
(Table~\ref{tab:confidence}), so that the certainty expressed in the
text is systematically aligned with the strength of the underlying
evidence.  This mapping may enable authors and reviewers to detect
confidence inflation---the failure mode described in Section~1---at the
drafting stage rather than during peer review.

\begin{table}[htbp]
\centering
\caption{Proposed confidence-to-language mapping.}
\label{tab:confidence}
\small
\begin{tabular}{@{}lll@{}}
\toprule
\textbf{Tier} & \textbf{Assign when\ldots}
  & \textbf{Prescribed language} \\
\midrule
Established & Multiple independent sources; no dispute
  & ``demonstrates,'' ``confirms'' \\
Supported   & 2--3 sources agree; open questions remain
  & ``indicates,'' ``evidence suggests'' \\
Emerging    & 1--2 sources; not yet replicated
  & ``may,'' ``preliminary evidence'' \\
Speculative & Logical inference; no direct data
  & ``warrants investigation'' \\
\bottomrule
\end{tabular}
\end{table}

% S3-4: Different types need different verification (ARGUMENT, P0, EMERGING)
The central argument of this paper is that these three types require
fundamentally different verification procedures, and that applying the
wrong procedure to the wrong type may produce systematic errors.  A
claim evaluated as an argument (checking its warrant rather than its
source) wastes effort.  More consequentially, an argument evaluated as a
claim (checking whether it has a source rather than whether its reasoning
is valid) may produce false failures---entries scored as unverified not
because they are wrong, but because the verification procedure was
inappropriate for the unit type.  This argument rests on the premise that
different epistemic contributions (facts, inferences, recommendations)
have different validity criteria; if this premise holds, then a unified
verification procedure will systematically misclassify some entries.
Preliminary evidence from the retrospective audits described in the
following section suggests that this type of misclassification may occur
in practice.


% ===============================================================
% SECTION 4 — PRELIMINARY EVIDENCE
% Registry entries: S4-1, S4-2, S4-3, S4-4
% ===============================================================
\section{Preliminary Evidence}

To explore whether typed verification reveals issues that conventional
approaches miss, we retrospectively applied the framework described in
Section~3 to three paper projects that had been developed with AI
writing assistance but without typed verification infrastructure.  These
audits should be understood as proof of concept rather than
validation---they are retrospective, conducted by the framework's
developers, and drawn from a single domain cluster.

% S4-1: Proposition audit (CLAIM, P0, EMERGING)
The most instructive case involved a proposition paper arguing that
feedback accuracy, not manikin realism, determines CPR training
effectiveness.  This paper had been developed with a conventional claim
registry that tracked twenty-one entries, achieving 76\% overall
coverage and 100\% coverage of critical claims.  A retrospective audit
reclassified all entries using the three-type model, revealing that
fifteen were indeed claims, but two were arguments and four were
propositions that had been misclassified as claims.

The consequences of this misclassification were most apparent for one
entry---the paper's central interpretive claim that invalid sensor
feedback undermines the primary benefit of technology-enhanced feedback
in CPR training.  Under claim-level verification, this entry was scored
at the lowest confidence tier (speculative, 0.25 on the project's
numeric scale) because it lacked a direct citation---it was classified as
``logical inference.''  However, this entry is not a claim; it is an
argument.  Its three premises---that sensors exhibit a 5.6\% calibration
shift when integrated into a manikin, that datasheet specifications
deviate 120--130 times from empirical calibration, and that most manikins
lack metrological characterisation---are each independently verified.
When evaluated using the Toulmin checklist, the entry achieves four of
five verification criteria: the claim is clearly stated, the grounds are
verified, the warrant is logically valid, and the qualifier is
appropriately calibrated.  Only counter-argument engagement is partial.
This represents what we term a \emph{false failure}: an entry scored as
the weakest in the registry not because its reasoning is unsound, but
because the wrong verification procedure was applied.  The distinction
matters for the manuscript: the writing guide mapped the speculative
score to heavily hedged language (``warrants investigation''), which is
excessively cautious for a logically sound argument whose premises are
verified.  Typed verification would have assigned an emerging confidence
tier, permitting language such as ``preliminary evidence suggests''---a
material improvement in argumentative force.

% S4-2: Technology audit (CLAIM, P1, EMERGING)
A second retrospective audit applied confidence tier analysis to a
technology paper under review at an IEEE instrumentation journal.  This
audit extracted twenty-two verifiable entries from the manuscript text
and assigned each a confidence tier based on the available evidence.  The
analysis identified six entries where the manuscript language was more
confident than the evidence warranted---for instance, using
``demonstrates'' for findings from a single validation study with a
limited sample size.

% S4-4: Cross-project applicability (ARGUMENT, P1, EMERGING)
Across all three retrospective audits---the proposition paper, the
technology paper, and an engineering characterisation paper targeting a
measurement science journal---the typed verification framework revealed
actionable issues in each project.  These issues included mistyped
entries producing false failures, over-confident language, missing
boundary conditions on propositions, and unregistered arguments in
discussion sections.  While the specific issues varied by paper type, the
common finding was that each audit identified problems that were not
visible without the typed verification infrastructure.

% S4-3: External validation (CLAIM, P1, EMERGING)
These preliminary findings are consistent with independent research on
structured verification for AI-assisted scholarly work.
\citet{peerarg2024} developed PeerArg, a framework that combines formal
argumentation structure with LLM-based review analysis; their results
suggest that structured representation of support and attack relations
improves review aggregation compared to end-to-end LLM approaches.
Similarly, \citet{gupta2024} found that prompting LLMs with Toulmin's
argument model significantly improves argument extraction compared to
generic prompts.  Both findings suggest that combining structured
frameworks with LLM capabilities may outperform LLM capabilities
alone---a principle consistent with the process-level verification
approach proposed here.

Several important limitations apply to these preliminary findings.  All
three audits were retrospective---the framework was applied after the
papers were drafted, not during writing.  All three projects share the
same author team and fall within a related domain cluster (medical
simulation and measurement science).  The audits were conducted by the
framework's developers, introducing potential confirmation bias.  These
limitations mean that the evidence should be interpreted as preliminary
indications rather than validation.  Whether typed verification
generalises to other authors, domains, and writing workflows remains an
open question.


% ===============================================================
% SECTION 5 — CALL TO ACTION
% Registry entries: S5-1, S5-2
% ===============================================================
\section{Call to Action}

The preceding sections have traced a chain of reasoning.  AI writing
assistance creates failure modes beyond citation hallucination, including
confidence inflation and argument-level degradation (Section~1).  The
most comprehensive verification infrastructure, EQUATOR's approximately
500 reporting guidelines, addresses empirical research exclusively;
non-empirical paper types---which include two of Gregor's five theory
types---have no equivalent verification framework (Section~2).
Established frameworks from argumentation theory (Toulmin) and
theory-building methodology (Whetten) may provide a foundation for typed
verification that distinguishes claims, arguments, and propositions
(Section~3).  Preliminary evidence from retrospective audits suggests
that this approach can reveal actionable issues, including false failures
caused by type misclassification (Section~4).

% S5-1: The proposition (PROPOSITION, P0, EMERGING)
On the basis of this reasoning, we propose that the scholarly community
develop verification infrastructure for non-empirical papers.  Such
infrastructure would extend the principle underlying
EQUATOR---that structured reporting improves research quality---to paper
types that currently have no systematic verification support.
Specifically, we propose that verification procedures be differentiated
by unit type: source verification for claims, Toulmin-based analysis for
arguments, and Whetten-based analysis for propositions.  Each type would
have associated checklists, confidence calibration criteria, and language
guidelines that authors can apply during manuscript preparation.

This proposal has specific boundary conditions.  It applies to academic
papers intended for peer-reviewed publication in fields where argument
quality constitutes a primary contribution.  It does not apply to
informal writing, journalism, creative writing, or very short opinion
pieces where the overhead of formal verification would exceed the
benefit.  Furthermore, the specific verification frameworks proposed
here---Toulmin's argument model and Whetten's theoretical contribution
framework---are grounded in Western academic argumentation traditions.
Their applicability to scholarly traditions with different rhetorical
conventions would need to be assessed rather than assumed.

One might object that journal peer review already provides adequate
argument verification for non-empirical papers---that reviewers routinely
evaluate logical coherence, boundary conditions, and counter-argument
engagement.  This observation has merit, and peer review remains
essential.  However, AI writing assistance may amplify the verification
challenge beyond what traditional peer review was designed to handle.
When AI tools can generate fluent, well-structured prose with minimal
effort, the volume and apparent quality of submitted manuscripts may
increase while the ratio of well-reasoned to superficially coherent
arguments decreases.  Reviewers already face growing workloads; adding
the burden of detecting AI-generated confidence inflation and implicit
warrant failures may exceed what the current system can absorb.
Process-level verification infrastructure would function as a complement
to peer review, not a replacement---catching argument-level issues before
submission so that reviewers can focus on substantive evaluation rather
than structural verification.

% S5-2: Next steps (CLAIM, P2, SPECULATIVE)
Several directions warrant investigation.  First, the typed verification
model proposed here has been developed and applied within medical
simulation and measurement science.  Discipline-specific verification
templates---adapting the three-type model to the conventions of different
fields---would be needed to assess the framework's generalisability.
Second, integration with AI writing tools could enable real-time
verification during the drafting process, using the structured checklists
and confidence calibration criteria as prompting infrastructure for
writing assistants.  Third, empirical validation comparing the quality of
papers written with and without typed verification infrastructure would
provide the evidence base needed to move the proposal from preliminary to
supported.

The EQUATOR Network has spent two decades building approximately 500
reporting guidelines for empirical health research.  These guidelines
have measurably improved the quality and transparency of published
research.  To our knowledge, essentially no equivalent guidelines exist
for non-empirical papers---theoretical contributions, design science
research, perspective articles, and methodological proposals.  As AI
writing tools make it easier to produce fluent academic prose, the gap
between what is published and what is rigorously verified may widen.  We
suggest it is time to begin building the verification infrastructure for
the rest.


% ===============================================================
% ACKNOWLEDGEMENTS
% ===============================================================
\section*{Acknowledgements}

[To be added.]


% ===============================================================
% REFERENCES
% ===============================================================
\bibliography{references}

\end{document}
